\uuid{I54c}
\titre{ Effet de la normalisation sur le gradient }
\niveau{L3}
\module{Analyse de données}
\chapitre{Réseaux de neurones}
\sousChapitre{Normalisation des données, descente de gradient}
\theme{Normalisation standard, gradient, vitesse de convergence}
\auteur{Maxime Nguyen}
\datecreate{2026-03-19}
\organisation{AMSCC}
\difficulte{2}

\contenu{
\texte{
  On dispose de deux exemples $x^{(1)} = 1000$ et $x^{(2)} = 2000$.
  La normalisation standard est effectuée avec une moyenne $\mu = 1500$ et un écart-type $\sigma = 500$.
}
\begin{enumerate}

\item
  \question{Calculer les valeurs normalisées $\tilde{x}^{(1)}$ et $\tilde{x}^{(2)}$.}
  \reponse{
    \[
      \tilde{x}^{(1)} = \frac{1000 - 1500}{500} = -1, \qquad
      \tilde{x}^{(2)} = \frac{2000 - 1500}{500} = 1.
    \]
  }

\item
  \question{Si on oublie de normaliser et qu'on initialise $w = 0{,}01$, comment le gradient $\frac{\partial J}{\partial w}$ se compare-t-il proportionnellement à celui obtenu avec les données normalisées ? Qu'est-ce que ça change en pratique ?}
  \reponse{
    Sans normalisation, le gradient est multiplié par un facteur de l'ordre de $x \sim 1000$--$2000$,
    soit environ $1000$ à $2000$ fois plus grand qu'avec des données normalisées (où $x \in \{-1, 1\}$).
    En pratique, cela impose d'utiliser un taux d'apprentissage très faible pour éviter la divergence,
    et ralentit considérablement la convergence.
  }

\end{enumerate}
}
